%  LaTeX support: latex@mdpi.com 
%  For support, please attach all files needed for compiling as well as the log file, and specify your operating system, LaTeX version, and LaTeX editor.

%=================================================================
\documentclass[minerals,article,submit,pdftex,moreauthors]{Definitions/mdpi} 

%=================================================================

% MDPI internal commands
\firstpage{1} 
\makeatletter 
\setcounter{page}{\@firstpage} 
\makeatother
\pubvolume{1}
\issuenum{1}
\articlenumber{0}
\pubyear{2023}
\copyrightyear{2023}
%\externaleditor{Academic Editor: Firstname Lastname}
\datereceived{19 February 2023} 
%\daterevised{} % Only for the journal Acoustics
\dateaccepted{} 
\datepublished{} 
\hreflink{https://doi.org/} % If needed use \linebreak
%\doinum{}

%=================================================================
% Add packages and commands here. The following packages are loaded in our class file: fontenc, inputenc, calc, indentfirst, fancyhdr, graphicx, epstopdf, lastpage, ifthen, lineno, float, amsmath, setspace, enumitem, mathpazo, booktabs, titlesec, etoolbox, tabto, xcolor, soul, multirow, microtype, tikz, totcount, changepage, attrib, upgreek, cleveref, amsthm, hyphenat, natbib, hyperref, footmisc, url, geometry, newfloat, caption

\graphicspath{{figures/}} % path to folder with figures
\usepackage{subcaption}
\usepackage{listings}
\usepackage{xcolor}
\usepackage{gensymb} % for \textdegree
\usepackage{tabularx}

\usepackage{listings}
\usepackage{xcolor}

\definecolor{deepblue}{rgb}{0,0,0.5}
\definecolor{deepred}{rgb}{0.6,0,0}
\definecolor{deepmagenta}{RGB}{195,12,214}
\definecolor{deepgreen}{RGB}{29,122,30}
\definecolor{mintgreen}{RGB}{192,249,192}
\definecolor{codepurple}{RGB}{204, 0, 153}
\definecolor{LightCyan}{rgb}{0.88,1,1}
\definecolor{GainsboroPurple}{RGB}{247,176,247}
\definecolor{LightSlateGrey}{RGB}{124,118,118}
%    
\lstdefinestyle{mystyle}{
    commentstyle=\color{LightSlateGrey},
    keywordstyle=\color{deepgreen}\bfseries,
    emphstyle=\color{deepred},
    numberstyle=\tiny\color{deepmagenta},
    stringstyle=\color{codepurple},
    basicstyle=\ttfamily\scriptsize,
    breakatwhitespace=false,         
    breaklines=true,                 
    captionpos=b,                    
    keepspaces=true,                 
    numbers=left,                    
    numbersep=5pt,                  
    showspaces=false,                
    showstringspaces=false,
    showtabs=false,                  
    tabsize=2
}
\lstset{style=mystyle}

%=================================================================
% Full title of the paper (Capitalized)
\Title{Title}

% MDPI internal command: Title for citation in the left column
\TitleCitation{Title}

% Author Orchid ID: enter ID or remove command
\newcommand{\orcidauthorA}{0000-0002-5759-1089} % Polina Add \orcidA{} behind the author's name
\newcommand{\orcidauthorB}{0000-0002-6461-1551} % Olivier Add \orcidB{} behind the author's name

% Authors, for the paper (add full first names)
\Author{Polina Lemenkova $^{1}$*\orcidA{} and Olivier Debeir $^{1}$\orcidB{}}

%\longauthorlist{yes}

% MDPI internal command: Authors, for metadata in PDF
\AuthorNames{Polina Lemenkova and Olivier Debeir}

% MDPI internal command: Authors, for citation in the left column
\AuthorCitation{Lemenkova, P.; Debeir, O.}
% If this is a Chicago style journal: Lastname, Firstname, Firstname Lastname, and Firstname Lastname.

% Affiliations / Addresses (Add [1] after \address if there is only one affiliation.)
\address{%
$^{1}$ \quad Laboratory of Image Synthesis and Analysis (LISA), École polytechnique de Bruxelles (Brussels Faculty of Engineering), Université Libre de Bruxelles (ULB). Building L, Campus du Solbosch, ULB -- LISA CP165/57, Avenue Franklin D. Roosevelt 50, Brussels 1050, Belgium.
}

% Contact information of the corresponding author
\corres{Correspondence: polina.lemenkova@ulb.be; Tel.: +32 471 86 04 59 (P.L.)}

% Current address and/or shared authorship
%\firstnote{Current address: Affiliation 3.} 

% Abstract (Do not insert blank lines, i.e. \\) 
\abstract{A single paragraph of about 200 words maximum. }

% Keywords
\keyword{keyword 1; keyword 2; keyword 3; keyword 4; keyword 5; keyword 6; keyword 7; keyword 8; keyword 9; keyword 10} 

% The fields PACS, MSC, and JEL may be left empty or commented out if not applicable
%\PACS{J0101}
%\MSC{}
%\JEL{}

%%%%%%%%%%%%%%%%%%%%%%%%%%%%%%%%%%%%%%%%%%
\begin{document}

\section{Introduction}

\begin{figure}[H]
\centering
	\includegraphics[width=14.0 cm]{Topo_CF.jpg}
	\caption{Topographic map of Central African Republic. Mapping software: Generic Mapping Tools (GMT) scripting toolset version 6.1.1, \cite{Wessel}. Data source: GEBCO/SRTM. Cartography source: authors.
	\label{fig01}}
\end{figure}

%%%%%%%%%%%%%%%%%%%%%%%%%%%%%%%%%%%%%%%%%%
\pagebreak
\section{Materials and Methods}

%----------- subsection ----------->
\subsection{Subsection}

\begin{lstlisting}[language=bash,caption=GMT script for computing the histogram equalisation of CAR's topographic grids,style=mystyle,label={lst01}]
exec bash
gmt grdcut ETOPO1_Ice_g_gmt4.grd -R14/28/2.5/11.5 -Gcf1_relief.nc
ps=HistCAR.ps
gmt makecpt -Crainbow -V -T212/1820 > t.cpt
gmt makecpt -Crainbow -T0/15/1 > c.cpt
# 1 upper left plot
gmt grdhisteq cf1_relief.nc -Gout.nc -C16
gmt grdimage cf1_relief.nc -I+a45+nt1 -Ct.cpt -JM3i -Y6i -K -P \
    -Bpxg4f1a2 -Bpyg4f2a2 -Bsxg2 -Bsyg2 -BWSNe > $ps
echo "24 10 Original" | gmt pstext -Rcf1_relief.nc -J -O -K -F+jBL+f12p -T -Gwhite@10 -Dj0.1i >> $ps
gmt pscoast -R -J -Ia/thinner,blue -Na -N1/thick,white -W0.1p -Df -O -K >> $ps
# 2 upper right plot
gmt grdimage out.nc -Cc.cpt -J -X3.5i -K -O \
    -Bpxg4f2a4 -Bpyg4f2a2 -Bsxg2 -Bsyg2 -BWSNe >> $ps
gmt pscoast -R -J -Ia/thinner,blue -Na -N1/thick,white -W0.1p -Df -O -K >> $ps
echo "24 10 Equalized" | gmt pstext -R -J -O -K -F+jBL+f12p -T -Gwhite@10 -Dj0.1i >> $ps
gmt psscale -Dx0i/-0.4i+jTC+w5i/0.15i+h+e+n -O -K -Ct.cpt -Bg200f10a200 -By+lm >> $ps
# 3 low left plot
gmt grdhisteq cf1_relief.nc -Gout.nc -N
gmt makecpt -Crainbow -T-3/3 > c.cpt
gmt grdimage out.nc -Cc.cpt -J -X-3.5i -Y-3.0i -K -O \
    -Bpxg4f2a4 -Bpyg4f2a2 -Bsxg2 -Bsyg2 -BWSNe >> $ps
gmt grdcontour cf1_relief.nc -R -J -C200 -A500+f7p,26,blue -Wthinnest,blue -O -K >> $ps
echo "23.5 10 Normalized" | gmt pstext -R -J -O -K -F+jBL+f12p -T \
    -UBL/-5p/-4.0c -Gwhite@10 -Dj0.1i >> $ps
gmt pscoast -R -J -Ia/thinner,blue -Na -N1/thick,white -W0.1p -Df -O -K >> $ps
# 4 low right plot
gmt grdhisteq cf1_relief.nc -Gout.nc -Q
gmt makecpt -Crainbow -T0/15 > q.cpt
gmt grdimage out.nc -Cq.cpt -J -X3.5i -K -O \
    -Bpxg4f2a4 -Bpyg4f2a2 -Bsxg2 -Bsyg2 -BWSNe >> $ps
gmt pscoast -R -J -Ia/thinner,blue -Na -N1/thick,white -W0.1p -Df -O -K >> $ps
echo "24 10 Quadratic" | gmt pstext -R -J -O -K -F+jBL+f12p -T -Gwhite@10 -Dj0.1i >> $ps
gmt psscale -Dx0i/-0.4i+w5i/0.15i+h+jTC+e+n -O -K -Cc.cpt -Bx1 -By+l"z@-n@-" >> $ps
gmt psscale -Dx0i/-1.0i+w5i/0.15i+h+jTC+e+n -O -K -Cq.cpt -Bx1 -By+l"z@-q@-" >> $ps
# GMT logo
gmt logo -Dx0.0/-4.5c+o-1.0c/0.2c+w2c -O -K >> $ps
# Subtitle
gmt pstext -R0/10/0/15 -JX10/10 -X-8.0 -Y6.5c -N -O \
-F+f14p,Palatino-Roman,black+jLB >> $ps << EOF
0.0 11.5 Central African Republic: histogram equalization of topographic grid:
1.5 10.7 ETOPO1 DEM Global Relief Model 1 arc min resolution
EOF
# Convert to image file by GhostScript
gmt psconvert HistCAR.ps -A0.5c -E720 -Tj -Z
\end{lstlisting}

%----------- subsection ----------->
\subsection{Subsection}

%----------- subsection ----------->
\subsection{Subsection}

%%%%%%%%%%%%%%%%%%%%%%%%%%%%%%%%%%%%%%%%%%
\section{Results}

\begin{figure}[H]
\makebox[1.0\linewidth][r]{%
%\centering
	\includegraphics[width=18.0 cm]{HistCAR.jpg}
}
	\caption{Histogram equalisation of the topography grids on CAR based on the ETOPO1 raster data. Mapping software: GMT v. 6.1.1, \cite{Wessel}. Data source: EMAG2. Cartography source: authors.
	\label{fig01}}
\end{figure}

\begin{figure}[H]
\centering
	\includegraphics[width=14.0 cm]{Magnet_CF.jpg}
	\caption{Map of magnetic anomaly in CAR. Mapping software: GMT v. 6.1.1, \cite{Wessel}. Data source: EMAG2. Cartography source: authors. Note: black and dark grey areas mean the "no data" areas where no magnetic survey was done.
	\label{fig01}}
\end{figure}

%%%%%%%%%%%%%%%%%%%%%%%%%%%%%%%%%%%%%%%%%%
\section{Discussion}

%%%%%%%%%%%%%%%%%%%%%%%%%%%%%%%%%%%%%%%%%%
\section{Conclusions}

%%%%%%%%%%%%%%%%%%%%%%%%%%%%%%%%%%%%%%%%%%
\authorcontributions{Supervision, conceptualization, methodology, software, resources, funding acquisition, and project administration, O.D.; writing---original draft preparation, methodology, software, data curation, visualization, formal analysis, validation, writing---review and editing, and investigation, P.L. All authors have read and agreed to the published version of the manuscript.}

\funding{The publication was funded by the Editorial Office of XXXX, Multidisciplinary Digital Publishing Institute (MDPI), by providing 100\% discount for the APC of this manuscript. This project was supported by the Federal Public Planning Service Science Policy or Belgian Science Policy Office, Federal Science Policy - BELSPO (B2/202/P2/SEISMOSTORM).}

\institutionalreview{Not applicable.}

\informedconsent{Not applicable.}

\dataavailability{Not applicable} 

\acknowledgments{The authors thank the reviewers for reading and review of this manuscript.}

\conflictsofinterest{The authors declare no conflict of interest.} 

\abbreviations{Abbreviations}{
The following abbreviations are used in this manuscript:\\

\noindent 
\begin{tabular}{@{}ll}
MDPI & Multidisciplinary Digital Publishing Institute\\
DOAJ & Directory of open access journals\\
LD & Linear dichroism
\end{tabular}
}

%%%%%%%%%%%%%%%%%%%%%%%%%%%%%%%%%%%%%%%%%%
%% Optional
\appendixtitles{no} % Leave argument "no" if all appendix headings stay EMPTY (then no dot is printed after "Appendix A"). If the appendix sections contain a heading then change the argument to "yes".
\appendixstart
\appendix
\section[\appendixname~\thesection]{}
\subsection[\appendixname~\thesubsection]{}
The appendix is an optional section that can contain details and data supplemental to the main text---for example, explanations of experimental details that would disrupt the flow of the main text but nonetheless remain crucial to understanding and reproducing the research shown; figures of replicates for experiments of which representative data are shown in the main text can be added here if brief, or as Supplementary Data. Mathematical proofs of results not central to the paper can be added as an appendix.

\begin{table}[H] 
\caption{This is a table caption.\label{tab5}}
\newcolumntype{C}{>{\centering\arraybackslash}X}
\begin{tabularx}{\textwidth}{CCC}
\toprule
\textbf{Title 1}	& \textbf{Title 2}	& \textbf{Title 3}\\
\midrule
Entry 1		& Data			& Data\\
Entry 2		& Data			& Data\\
\bottomrule
\end{tabularx}
\end{table}

\section[\appendixname~\thesection]{}
All appendix sections must be cited in the main text. In the appendices, Figures, Tables, etc. should be labeled, starting with ``A''---e.g., Figure A1, Figure A2, etc.

%%%%%%%%%%%%%%%%%%%%%%%%%%%%%%%%%%%%%%%%%%
\begin{adjustwidth}{-\extralength}{0cm}
%\printendnotes[custom] % Un-comment to print a list of endnotes

\reftitle{References}
%\nocite{*}

%=====================================
% References, variant A: external bibliography
%=====================================
\bibliography{Bibliography.bib}

%%%%%%%%%%%%%%%%%%%%%%%%%%%%%%%%%%%%%%%%%%
\end{adjustwidth}
\end{document}
